% !TEX program = pdflatex
% In-depth explanatory companion to main2.tex.
% This document explains the project proof mechanics; it does not claim P != NP.

\documentclass[11pt]{article}
\usepackage[T1]{fontenc}
\usepackage[utf8]{inputenc}
\usepackage[margin=0.86in]{geometry}
\usepackage{amsmath,amssymb,mathtools,amsthm}
\usepackage{booktabs,longtable,tabularx,array}
\usepackage[table]{xcolor}
\usepackage{enumitem}
\usepackage[most]{tcolorbox}
\usepackage{hyperref}
\usepackage{microtype}
\usepackage{fancyhdr}
\usepackage{lastpage}

\definecolor{proved}{HTML}{1B7F3A}
\definecolor{openstatus}{HTML}{B15C00}
\definecolor{retired}{HTML}{A12A2A}
\definecolor{candidate}{HTML}{4B5FA5}
\definecolor{softblue}{HTML}{EEF4FA}
\definecolor{softgreen}{HTML}{EDF8F0}
\definecolor{softorange}{HTML}{FFF4E8}
\definecolor{softred}{HTML}{FCEEEE}
\definecolor{softgray}{HTML}{F5F6F7}

\hypersetup{
 colorlinks=true,
 linkcolor=candidate,
 urlcolor=candidate,
 citecolor=candidate,
 pdftitle={P vs NP Research Program --- In-Depth Explanation of Core Proof Mechanics},
 pdfauthor={P vs NP Research Project}
}

\pagestyle{fancy}
\fancyhf{}
\lhead{\small P vs NP Research Program}
\rhead{\small In-Depth Proof Mechanics}
\cfoot{\small \thepage\ /\ \pageref{LastPage}}
\setlength{\headheight}{14pt}

\newcommand{\same}{\textsc{same\_realization}}
\newcommand{\poly}{\operatorname{poly}}
\newcommand{\statusproved}{\textcolor{proved}{\textbf{PROVED/EXACT}}}
\newcommand{\statusopen}{\textcolor{openstatus}{\textbf{OPEN}}}
\newcommand{\statuscandidate}{\textcolor{candidate}{\textbf{CANDIDATE}}}
\newcommand{\statusretired}{\textcolor{retired}{\textbf{FALSIFIED/RETIRED}}}

\newtcolorbox{statusbox}{
 colback=softorange,colframe=openstatus,boxrule=0.8pt,arc=2pt,
 left=7pt,right=7pt,top=6pt,bottom=6pt
}
\newtcolorbox{provedbox}{
 colback=softgreen,colframe=proved,boxrule=0.7pt,arc=2pt,
 left=7pt,right=7pt,top=6pt,bottom=6pt
}
\newtcolorbox{candidatebox}{
 colback=softblue,colframe=candidate,boxrule=0.7pt,arc=2pt,
 left=7pt,right=7pt,top=6pt,bottom=6pt
}
\newtcolorbox{warningbox}{
 colback=softred,colframe=retired,boxrule=0.7pt,arc=2pt,
 left=7pt,right=7pt,top=6pt,bottom=6pt
}
\newtcolorbox{explainbox}{
 colback=softgray,colframe=black!45,boxrule=0.55pt,arc=2pt,
 left=7pt,right=7pt,top=6pt,bottom=6pt
}

\theoremstyle{plain}
\newtheorem{theorem}{Theorem}[section]
\newtheorem{lemma}[theorem]{Lemma}
\newtheorem{corollary}[theorem]{Corollary}
\theoremstyle{definition}
\newtheorem{definition}[theorem]{Definition}
\theoremstyle{remark}
\newtheorem{remark}[theorem]{Remark}

\title{\bfseries P vs NP Research Program\\[0.25em]
\Large In-Depth Explanation of the Core Proof Mechanics}
\author{Companion exposition to \texttt{main2.tex}}
\date{2026-10-05}

\begin{document}
\maketitle

\begin{statusbox}
\textbf{Global status and purpose.}
\[
\boxed{\text{$P$ versus $NP$ remains OPEN.}}
\]
This document is an explanatory companion to \texttt{main2.tex}.  It is meant
to make the proof architecture readable from first principles: what each
coordinate means, why each theorem was introduced, which implications are
actually proved, which ideas were falsified, and what would still have to be
proved for an unconditional separation.

Unless explicitly labeled otherwise, the mathematical claims below are
explanations of results and status statements already present in the project
corpus.  This file does not promote a restricted theorem, finite computation,
or conditional implication into an unrestricted complexity lower bound.
\end{statusbox}

\tableofcontents
\bigskip

\section{Executive overview}

The research program has gradually converged on a fairly sharp diagnosis of
the P-vs-NP difficulty.

It is easy to write down many quantities that become large on a hard-looking
3-SAT formula: number of satisfying assignments, number of possible histories,
number of clauses, number of provenance labels, rank of a matrix, width of a
decision representation, size of an explicit strategy tree, or number of
distinguishable future continuations.  None of these is automatically a time
lower bound for an arbitrary algorithm.

The reason is representation freedom.  An algorithm can summarize, reuse,
recompute, reread, share a DAG node, switch to an algebraic basis, invoke a
dynamic program on a separator, or use a completely different representation
from the one the proof is counting.

The current program therefore separates the intended proof into three layers.

\begin{enumerate}[leftmargin=1.6em]
\item \textbf{Semantic layer.}
Define exactly which distinctions remain relevant to the future consumer.

\item \textbf{Source/accounting layer.}
Track where the information supporting those distinctions lives, how it is
represented, and whether it can be recovered or recomputed.

\item \textbf{Cost layer.}
Prove that every ordinary way of maintaining or reacquiring enough of those
distinctions costs superpolynomial resources.
\end{enumerate}

The first two layers are now comparatively mature.  The third is still the
central open problem.

\begin{candidatebox}
\textbf{Current intended separation spine.}
\[
\boxed{
\begin{array}{c}
\text{hypothetical polynomial-time exact SAT decider }A\\
\Downarrow\\
\text{SAT-specific source-to-reduction extraction (OPEN)}\\
\Downarrow\\
n-O(\log n)\text{ valid future-relevant reduction}\\
\Downarrow\\
\text{SRAS route completeness (proved in the declared interface model)}\\
\Downarrow\\
\text{retain/localize}\ \vee\ \text{reacquire}\ \vee\ \text{external}\\
\Downarrow\\
\text{ordinary retained/reacquired branch costs }n^{\omega(1)}\text{ (OPEN)}\\
\Downarrow\\
P\ne NP.
\end{array}}
\]
\end{candidatebox}

The exact-3SAT work runs in parallel.  It is not itself the unrestricted proof.
Its purpose is to expose transformations that are genuinely cheap so that the
abstract lower-bound argument is not allowed to charge them as hard.

\section{Three distinctions that govern the whole project}

\subsection{Semantic multiplicity is not runtime}

Suppose a formula has \(2^n\) assignments.  That does not imply an algorithm
must examine \(2^n\) objects.

Gaussian elimination on XOR-SAT is the canonical control.  An affine system
can have exponentially many assignments and strong global dependencies, yet a
compact basis decides the problem in polynomial time.

So the implication
\[
\text{many semantic possibilities}
\Longrightarrow
\text{many computational steps}
\]
is false without an additional theorem.

\subsection{Existence is not acquisition}

Suppose there exists a short witness \(Q\) that would answer all relevant
future queries.  A proof still has to explain how the computation obtains
\(Q\).

Conversely, if \(Q\) is a polynomial-time function of the input, then
discarding it and recomputing it later may be perfectly cheap.

This is why the project repeatedly separates:
\[
\boxed{\text{semantic existence}\ne\text{computational acquisition}.}
\]

\subsection{Representation size is not source size}

A tiny object can depend on the entire input.

A checksum is the simplest analogy:
\[
Q(F)\in\{0,1\}^{O(\log n)}
\]
may be a very small representation even though
\[
|\operatorname{Src}(Q)|=\Theta(|F|).
\]

Thus:
\[
\boxed{\text{small representation}\not\Rightarrow\text{small source}.}
\]

CSE needs actual source contraction, not merely a short message.

\section{The Boolean reduction threshold}

The first calculation in \texttt{main2.tex} is intentionally elementary,
because it calibrates the amount of reduction a polynomial endpoint would
require.

Start with a conceptual space of size
\[
2^n.
\]

Suppose a valid reduction eliminates a factor \(2^x\).  The residual scale is
\[
2^{n-x}.
\]

If we want this residual scale to be polynomial,
\[
2^{n-x}=n^d,
\]
then taking logarithms gives
\[
\boxed{x=n-d\log_2 n.}
\]

Equivalently, define the residual exponent
\[
D(n)=n-x(n).
\]

Then:
\[
\boxed{
\begin{array}{rcl}
D(n)&=&O(\log n)\quad\Rightarrow\quad\text{polynomial residual},\\[0.3em]
D(n)&=&\omega(\log n)\quad\Rightarrow\quad\text{superpolynomial residual},\\[0.3em]
D(n)&=&\Omega(n)\quad\Rightarrow\quad\text{exponential residual}.
\end{array}}
\]

\begin{explainbox}
This theorem does \emph{not} say every polynomial SAT algorithm explicitly
removes \(n-O(\log n)\) assignment bits.  It says that if we independently
define a Boolean reduction measure whose residual interpretation is
\(2^{D(n)}\), then polynomial residual scale requires \(D(n)=O(\log n)\).

The missing step is precisely the arbitrary-algorithm theorem connecting an
unknown SAT computation to such a reduction object.
\end{explainbox}

\section{Continuation equivalence: what may safely be forgotten}

\subsection{Why raw history counting is too strong}

A computation can have many different histories that are irrelevant to every
future action.

If two histories differ only in dead code, an unused random bit, or the order
in which an already-computed fact was obtained, forcing the machine to
distinguish them would be artificial.

The project therefore defines equivalence relative to a declared future
consumer universe.

\begin{definition}[Continuation equivalence]
Let \(\mathcal H\) be a family of histories or residual situations and \(U\) a
family of legal future continuations.  Define
\[
h\equiv_U h'
\iff
\forall u\in U,\quad
\Phi(h,u)=\Phi(h',u).
\]
Let
\[
\kappa_U(\mathcal H)
=
|\mathcal H/\!\equiv_U|.
\]
\end{definition}

Only histories in different continuation classes are semantically required to
remain distinguishable for the declared consumer.

\subsection{The safe-compression lemma}

Suppose a summary
\[
\sigma:\mathcal H\to Z
\]
supports a deterministic decoder satisfying
\[
D(\sigma(h),u)=\Phi(h,u)
\]
for all legal \(h,u\).

If
\[
h\not\equiv_Uh',
\]
then
\[
\sigma(h)\ne\sigma(h').
\]

Therefore:
\[
\boxed{
|\operatorname{range}\sigma|
\ge
\kappa_U(\mathcal H).
}
\]

This is a clean pigeonhole theorem.  It is also exactly where many false
lower-bound arguments must stop.

\begin{warningbox}
A lower bound on the number of required summary states is not automatically a
time lower bound.  The state index may have a compact representation, may be
computed algebraically, or may already be present in the input representation.
\end{warningbox}

\section{SRAS: the Source-Reduction Accounting System}

\subsection{What SRAS is trying to accomplish}

SRAS is not itself a solver and not itself a lower bound.  It is a ledger.

Whenever the proof says that a future-relevant distinction disappeared, SRAS
asks:
\begin{itemize}[leftmargin=1.4em]
\item Was it truly shown irrelevant?
\item Was it merged with an equivalent state?
\item Was it encoded into a different representation?
\item Was it localized to a smaller source?
\item Can the computation reread or recompute it?
\item Was it supplied externally?
\end{itemize}

The purpose is to prevent a proof from silently converting information loss
into computational cost.

\subsection{The full state tuple}

The current SRAS state is
\[
\mathcal A_t=
(X,K,Y,Z,H,J,P,B,U,\Phi_U,V,\Psi_V,R,E,C,Q,G,L,W,\mathsf A,
\Omega_D,\Omega_G,\mathrm{Fate}).
\]

The coordinates are intentionally redundant because they separate concepts
that were previously conflated.

\begin{longtable}{p{0.11\textwidth} p{0.79\textwidth}}
\toprule
\textbf{Coordinate} & \textbf{Interpretation}\\
\midrule
\endhead
\(X\) & Original source identity or source object.\\
\(K\) & Stable identity of the clause, residual object, or canonical object.\\
\(Y,Z\) & Current semantic values or relations.\\
\(H\) & Realized derivation/audit history.\\
\(J\) & Source or causal lineage.\\
\(P\) & Provenance tags, prime tags, hashes, or other identity labels.\\
\(B\) & Boundary/separator/interface through which the object is consumed.\\
\(U\) & Declared continuation universe.\\
\(\Phi_U\) & Required future semantic behavior.\\
\(V,\Psi_V\) & Provenance-sensitive observations and their required behavior.\\
\(R\) & Representation currently retained.\\
\(E\) & Enabled operations and legal escape mechanisms.\\
\(C\) & Certificates already established.\\
\(Q\) & Reacquisition/reconstruction method.\\
\(G\) & Combination, reuse, and sharing rule.\\
\(L\) & Loss/quotient event.\\
\(W\) & Recoverability of discarded information.\\
\(\mathsf A\) & Access path to the relevant source/representation.\\
\(\Omega_D\) & Data authority: input, memory, randomness, advice, etc.\\
\(\Omega_G\) & Operation authority: ordinary implementation, subroutine, oracle.\\
Fate & Final disposition of each future-relevant contribution.\\
\bottomrule
\end{longtable}

\subsection{Lineage hierarchy}

For a derived object \(q\), the project uses
\[
\boxed{
J_{\rm live}^{U}(q)
\subseteq
J_{\rm sem}(q)
\subseteq
J_{\rm syn}(q)
\subseteq
J_{\rm origin}(q).
}
\]

These sets answer different questions.

\begin{itemize}[leftmargin=1.4em]
\item \(J_{\rm origin}\): what ever causally contributed?
\item \(J_{\rm syn}\): what does the current syntax refer to?
\item \(J_{\rm sem}\): what actually affects the current meaning?
\item \(J_{\rm live}^{U}\): what can still affect a legal future consumer?
\end{itemize}

This is why a large original dependency set need not remain expensive.  Many
dependencies can become irrelevant or safely compressed.

\subsection{Fate labels}

Each future-relevant contribution eventually receives one of the following
dispositions.

\begin{description}[leftmargin=1.8em,style=nextline]
\item[LIVE] A legal future consumer still distinguishes it.
\item[ENCODED] It remains available through another representation.
\item[LOCALIZED] Its live dependence has been confined to a smaller source.
\item[QUOTIENTED] Distinct histories were safely merged.
\item[COMBINED] Several contributions are represented jointly.
\item[IRRELEVANT] Future consumers provably do not need it.
\item[TERMINAL] No future consumer remains.
\item[EXTERNAL] It is supplied by an external authority.
\end{description}

\subsection{Valid continuation reduction}

If \(\kappa_U(S)\) is the number of live continuation classes in state \(S\),
define
\[
\boxed{
K_{s\to t}
=
\log_2
\frac{\kappa_U(S_s)}{\kappa_U(S_t)}.
}
\]

This is the project's preferred reduction quantity because it is defined from
future semantics, not from the observed runtime of the candidate algorithm.

\subsection{The acquisition curve}

For a formula \(F\), consumer \(U\), and computation model \(M\), define
\[
\boxed{
\mathcal C_{F,U,M}(k)
=
\min_{\mathcal D}
\left\{
C(\mathcal D):
K_U(\mathcal D)\ge k
\right\}.
}
\]

The desired lower-bound target is
\[
\boxed{
\mathcal C_{F_n,U_n,M}
\bigl(n-O(\log n)\bigr)
=
n^{\omega(1)}.
}
\]

This is a much stronger statement than saying the formula has many semantic
states.  It says every legal ordinary computation DAG that acquires enough
future-relevant reduction must pay superpolynomial cost.

That theorem is \statusopen.

\section{SRAS route completeness}

\subsection{The causal conservation argument}

One part of SRAS is now an exact theorem in its declared deterministic
interface model.

At time \(t\), suppose all future behavior can depend only on:
\[
R_t=\text{retained state},
\]
\[
A_t=\text{declared accessible/re-readable source},
\]
and future external inputs \(E_{>t}\).

If two histories have equal \(R_t\), equal accessible source answers, and
receive equal external inputs, then deterministic transition induction forces
all future states and outputs to be equal.

Therefore, if two histories require different future responses, the
distinguishing contribution must be somewhere:
\[
\boxed{
\text{retained/encoded}
\vee
\text{recoverable from accessible source}
\vee
\text{future external information}.
}
\]

Quotienting is legal only if the histories are already continuation-equivalent.

\begin{provedbox}
This closes a bookkeeping question: there is no mysterious fourth place for a
future distinction to live in the declared deterministic interface.

It does \emph{not} prove that retaining or reacquiring the distinction is
expensive.
\end{provedbox}

\section{CSE: why extraction remains the main arrow}

\subsection{Trace extraction is not enough}

If a deterministic algorithm runs in polynomial time, then its explicit
execution trace is polynomial-size.  Any polynomial-time transformation of
that trace is also polynomial-time.

Thus:
\[
\boxed{
\text{polynomial trace}
+
\text{polynomial audit}
\not\Rightarrow
\text{superpolynomial resource}.
}
\]

This is TRACECLOSE1.

\subsection{A fixed consumer is available}

The project once worried that a proof might invent a counterfactual consumer
not actually forced by SAT correctness.

EXTCONSUMER1 eliminates that ambiguity by using:
\[
\boxed{
W_{\rm EXT}(F,\rho)
=
\operatorname{SAT}(F\upharpoonright\rho).
}
\]

Every exact SAT decider \(A\) provides:
\[
E_A(F,\rho)
=
A(F\upharpoonright\rho).
\]

So the universal extension consumer is solver-independent.

But this still does not contract source scale.  A one-variable restriction
typically maps
\[
N\mapsto N-O(1),
\]
not
\[
N\mapsto N^b,\qquad b<1.
\]

\subsection{The CSE contract}

A usable Contractive Source Extraction theorem must output recursively
consumable objects \(Q_1,\ldots,Q_m\) with:
\[
m\le N^{a+o(1)},
\]
\[
|J_{\rm live}^{U}(Q_i)|
\le
N^{b+o(1)},
\qquad b<1,
\]
and
\[
C_{\rm acquire}
+
C_{\rm combine}
\le
N^{c+o(1)}.
\]

It must also satisfy:
\begin{itemize}[leftmargin=1.4em]
\item exact correctness;
\item recursive closure;
\item no hidden current-scale SAT call;
\item wrapper covariance;
\item source matching;
\item explicit acquisition/reconstruction cost;
\item no free canonicalization.
\end{itemize}

This is still \statusopen.

\section{What recursive contraction would buy}

Suppose a valid extractor/shrinker did exist.

If local work costs
\[
n^{r+o(1)},
\]
there are at most
\[
n^{\eta+o(1)}
\]
recursive calls, and child size is
\[
n^{\theta+o(1)},
\qquad 0\le\theta<1,
\]
then a polynomial exponent \(C\) satisfies
\[
C
=
\max\left\{
r,\eta+\theta C
\right\},
\]
so
\[
\boxed{
C
=
\max\left\{
r,
\frac{\eta}{1-\theta}
\right\}.
}
\]

This is RECURCLOSE1.

The importance of this theorem is conceptual.  It removes the unknown starting
polynomial exponent if the shrinker has genuine fractional contraction.

By contrast, if a supposedly recursive transform makes a call at current
scale, then SELECTCHARGE1 shows the unknown exponent simply reappears:
\[
n^{u+\delta K},
\]
with \(\delta=1\) restoring \(K\).

\section{Why provenance and replay did not solve the problem}

\subsection{What provenance does solve}

Provenance can prevent the proof from replacing the actual computed object by
an unrelated semantically equivalent object.

This is the purpose of \same.

For example, two histories may satisfy:
\[
\operatorname{Sem}(h_1)
=
\operatorname{Sem}(h_2)
\]
while having different sources or derivations.

A provenance-sensitive consumer may therefore distinguish them.

\subsection{What provenance does not solve}

Perfect provenance can still show that every useful intermediate object depends
on the entire parent input.

For a compact object \(Q\):
\[
|Q|=O(\log n)
\]
can coexist with
\[
|\operatorname{Src}(Q)|=\Theta(N).
\]

Provenance labels identify that fact; they do not make the computation
expensive.

\subsection{Replay attacks}

Four attacks repeatedly defeated provenance-only lower bounds.

\begin{enumerate}[leftmargin=1.6em]
\item \textbf{Rerun.}
If \(Q(F)\) is polynomial-time, simply recompute it.

\item \textbf{Reread.}
If the static input remains available, information discarded from memory may
be reacquired from the input.

\item \textbf{Transcript regeneration.}
A polynomial functional of a polynomial trace stays polynomial.

\item \textbf{Quotient retention.}
A long history may compose into one polynomial-size summary.
\end{enumerate}

Therefore:
\[
\boxed{
\text{unique provenance}
\not\Rightarrow
\text{expensive acquisition}.
}
\]

\section{Random challenges: what survived}

The random-digit-string experiments produced one exact lesson and one no-go.

If \(\tau\) is a stopping time for IID fair bits, then:
\[
(R_{\tau+1},\ldots,R_{\tau+b})
\mid
\mathcal F_\tau
\sim
U(\{0,1\}^b).
\]

So a challenge sampled after \(\tau\) is genuinely fresh relative to the prior
causal state.

But freshness only prevents challenge-specific advance knowledge.  If the
computation retains a universal evaluator
\[
E(R,h,u),
\]
then it can answer after \(u\) is revealed.

Thus:
\[
\boxed{
\text{freshness}+\same
\not\Rightarrow
\text{expensive response}.
}
\]

The meaningful target becomes the complexity of the universal evaluator.

\section{Bounded-width closure and why hidden local inference is cheap}

A width-\(w\) CNF over \(n\) variables has at most
\[
N_w(n)
=
\sum_{i=0}^{w}
2^i\binom ni
\]
non-tautological clauses of width at most \(w\).

For constant \(w\),
\[
N_w(n)=n^{O(w)}.
\]

Therefore exhaustive width-\(w\) resolution saturation is polynomial in \(n\).

This has an important methodological consequence.

If the only reason a useful relation seems hard is that it takes several
constant-width inference steps to expose, that is not enough for an
unrestricted lower bound.  A global polynomial preprocessing pass can
enumerate the entire constant-width closure.

This is why mediator-hidden binary relations, guarded 2-SAT constraints, and
similar structures must be credited as cheap once discovered.

\section{Relativization: from a conceptual location to ASCS}

\subsection{Why generic SRAS arguments relativize}

Continuation equivalence, state counting, and route completeness do not care
whether one transition is computed by ordinary code or by an oracle.

Therefore these tools alone cannot resolve P versus NP.

Baker--Gill--Solovay imply that some load-bearing step must fail under arbitrary
oracle replacement.

\subsection{Source transparency}

The project candidate is ordinary-source transparency.

An ordinary machine step has an implementation that can in principle be
expanded and charged.

An oracle step
\[
y=O(x)
\]
is opaque by definition.

This gives a place where a proof may stop under relativization without
declaring oracle calls expensive by fiat.

\subsection{Absolute Source-Complexity Separator}

The Google Drive source audit contains a more concrete construction.

Fix an ordinary universal machine \(U\), public conditioning \(c\), output
length \(m\), description budget \(B\), and time bound \(\tau\).  Define:
\[
\mathrm{Low}_U(c;m,B,\tau)
=
\left\{
y\in\{0,1\}^m:
\exists q,\ |q|\le B,\ U(q,c)\text{ outputs }y\text{ within }\tau
\right\}.
\]

Then:
\[
D_{\rm ASCS}(y)=1
\iff
y\notin\mathrm{Low}_U(c;m,B,\tau).
\]

There are fewer than \(2^{B+1}\) descriptions of length at most \(B\), hence:
\[
\boxed{
|\mathrm{Low}_U|<2^{B+1}.
}
\]

If \(B\le m-2\), then ASCS accepts at least half of the \(m\)-bit cube.

If a same-realization RMV range \(R_\pi\) is contained in
\(\mathrm{Low}_U\), ASCS misses every point of that range.

\subsection{Why ASCS genuinely breaks oracle replacement}

Choose a string \(y_m\) outside the ordinary-source set
\(\mathrm{Low}_U\).

An oracle can simply encode the bits of \(y_m\).  A short oracle program then
outputs \(y_m\) quickly, even though no corresponding short ordinary program
does so within the prescribed time.

Thus the ordinary-source containment premise can fail after:
\[
U\to U^O.
\]

This is a genuine oracle-break diagnostic.

\begin{provedbox}
ASCS gives an explicit example of a source-semantic statement that need not
survive oracle replacement.

It still does not prove a SAT lower bound.
\end{provedbox}

\section{FSC: the verification bottleneck inside ASCS}

To use ASCS as the small co-nondeterministic distinguisher required by the RMV
reconstruction, the complement relation must have a sufficiently small
verification circuit.

The desired Fast Source Certificate says, roughly:

for every relevant output \(y\), there is a short ordinary source description
and an exact derivation/evaluation certificate whose deterministic verifier
fits in the target size \(m\).

In direct form, seek \(V_{\rm src}\) of size at most \(m\) satisfying:
\[
y\in R_\pi
\iff
\exists w\ V_{\rm src}(c,y,w)=1.
\]

Then:
\[
D(y)
=
\forall w\,[V_{\rm src}(c,y,w)=0]
\]
is the desired co-nondeterministic separator.

The problem is that direct replay may cost the original source production
time \(T\), which can be much larger than \(m\).

Probabilistic checking does not immediately fix this because the required
relation is exact.  Algebraic verification raises a separate algebrization
audit.

Thus FSC remains \statusopen.

\begin{explainbox}
The ASCS/FSC line isolates a useful tension:

\[
\boxed{
\begin{array}{c}
\text{easy extensional image test}
\Rightarrow
\text{compact but typically relativizing},\\
\text{absolute source complexity}
\Rightarrow
\text{can fail oracle replacement},\\
\text{exact tiny verification of that source relation}
\Rightarrow
\text{currently open}.
\end{array}}
\]
\end{explainbox}

\section{Algebrization and natural-proofs discipline}

Nonrelativizing is not the same as nonalgebrizing.

Low-degree extensions, Reed--Muller encodings, polynomial identities, and
algebraic reconstruction live close to the methods that the
Aaronson--Wigderson barrier was designed to audit.

Therefore any source theorem using algebraic encodings must identify which
load-bearing step fails even when an algebraic oracle extension is allowed.

Likewise, an efficiently recognizable large property of truth tables would
need a Razborov--Rudich natural-proofs audit if it were used as a circuit
lower-bound property.

The project currently has no proved final nonalgebrizing hinge.

\section{The exact-3SAT laboratory}

The exact-3SAT branch serves as both an algorithmic project and a red-team
environment for SRAS.

If the exact solver can expose a relation cheaply, SRAS must treat that
relation as cheap.

The current laboratory is built around Liu's chain framework.

\subsection{Type-1 chain}

A positive type-1 chain is one positive 3-clause:
\[
x_1\vee x_2\vee x_3.
\]

Its satisfying local state space is:
\[
D
=
\{0,1\}^3\setminus\{000\},
\qquad
|D|=7.
\]

Its Liu characteristic is:
\[
\boxed{
\lambda_1=\frac37.
}
\]

The current internal baseline for the type-1 balance is approximately:
\[
c_0\approx1.3279235672.
\]

This constant belongs to the restricted exact-algorithm analysis.  It is not
a lower bound for arbitrary SAT algorithms.

\section{Mediator elimination}

Suppose an outside variable \(y\) appears as:
\[
G_y
=
\bigwedge_{i\in P}(A_i\vee y)
\wedge
\bigwedge_{j\in N}(B_j\vee\neg y).
\]

Eliminating \(y\) gives:
\[
\exists y\,G_y
\iff
\bigwedge_{i\in P}
\bigwedge_{j\in N}
(A_i\vee B_j).
\]

So a constant-degree outside mediator can be removed by polynomially many
constant-width resolvents.

In SRAS language:
\begin{itemize}[leftmargin=1.4em]
\item the mediator can become irrelevant;
\item the induced relation is encoded in the resolvent;
\item the proof is not allowed to charge future rediscovery of a relation that
ordinary preprocessing can materialize cheaply.
\end{itemize}

\section{Why the clean-core syntax can still be NP-hard}

The project constructs an externally mediated clean-core class in which:
\begin{itemize}[leftmargin=1.4em]
\item base clauses are disjoint positive type-1 chains;
\item outside mediators have degree two and both polarities;
\item mediator pairs encode binary consistency conditions;
\item direct non-base chain-only clauses are absent.
\end{itemize}

A linear reduction from arbitrary 3-SAT shows that this schema remains
NP-complete.

This is important because it rules out the naive inference:
\[
\text{low mediator degree}
\Rightarrow
\text{easy residual problem}.
\]

At the same time, mediator elimination exposes useful finite-domain
correlations.  So NP-hard expressiveness and exploitable exact structure can
coexist.

\section{Exact relation characteristics}

For the seven-state chain domain define:
\[
K_{u,v}=2^{-d(u,v)},
\qquad
N=K^{-1},
\]
\[
z=N\mathbf1,
\qquad
t=\mathbf1^\top z=\frac73.
\]

For a local falsifying set \(B\subseteq D\), define:
\[
a(B)
=
z_B^\top
\bigl(N[B,B]\bigr)^{-1}
z_B.
\]

For a resolved clause \(C\) touching \(r\) chains with local falsifying sets
\(B_1,\ldots,B_r\):
\[
\boxed{
\lambda(C)
=
\left[
\left(\frac73\right)^r
-
\prod_{i=1}^{r}a(B_i)
\right]^{-1}.
}
\]

This gives one closed form for all single resolved clauses in the current
rank-\(2/3/4\) schema.

The exhaustive worst cases are:
\[
\lambda_{\min,2}=\frac{81}{440},
\]
\[
\lambda_{\min,3}=\frac{729}{9236},
\]
\[
\lambda_{\min,4}=\frac{6561}{193856}.
\]

Relative to independent chains, the conservative DLS multipliers are:
\[
g_2=\frac{440}{441},
\qquad
g_3=\frac{9236}{9261},
\qquad
g_4=\frac{193856}{194481}.
\]

All are strictly below \(1\).

Thus every non-tautological resolved relation gives some positive local DLS
credit.

\section{Hypermatching: exploiting spread-out relations}

Create a rank-\(\le4\) hypergraph whose vertices are clean chains and whose
hyperedges are resolved relations.

If there are \(M\) hyperedges and each chain occurs in at most \(\Delta\)
hyperedges, a maximal hypermatching has at least:
\[
\frac{M}{4\Delta}
\]
edges and covers at least:
\[
\frac{M}{2\Delta}
\]
chains.

After removing inactive chains, ACTIVE-DENSITY1 gives:
\[
M\ge\frac{\nu_{\rm act}}4.
\]

Hence:
\[
\boxed{
f_{\rm covered,act}
\ge
\frac1{8\Delta}.
}
\]

When \(\Delta\) is bounded by a constant, this gives a constant positive
coverage fraction.

The difficulty had been the concentrated regime where \(\Delta\) becomes very
large.

\section{Link hypergraph and the former concentration escape}

Fix a chain \(T\) of large incidence.

Every resolved clause in the declared model contains at most two literals from
\(T\).

There are only 18 possible nonempty one-/two-literal local patterns on \(T\).

Pigeonhole by local pattern, remove \(T\), and consider the outside supports.
These supports have size at most three.

If the outside-support hypergraph has a large matching, we obtain many
pairwise-disjoint petals immediately.

If not, the union of a maximal matching forms a small vertex cover.  Pigeonhole
again on one cover vertex and its local pattern.  This enlarges the common
center.

Because the original rank is at most four, this process has bounded depth.

This is LINK-DICHOTOMY1.

\section{Closing the finite high-incidence pattern audit}

The 18 patterns collapse by coordinate symmetry to five sign types:
\[
1+,\quad 1-,\quad 2++,\quad 2+-,\quad 2--.
\]

For pairwise-disjoint exposed petals, use the conservative factor:
\[
g=\frac{440}{441}.
\]

For each center type the branch ratio takes the form:
\[
R_Q(m)=a_Q+b_Qg^m.
\]

The exact table is:

\begin{center}
\begin{tabular}{c|c|c|c}
\toprule
Center type & \(a_Q\) & \(b_Q\) & First \(m\) with \(R_Q(m)<1\)\\
\midrule
\(1+\) & \(16/21\) & \(5/7\) & 484\\
\(1-\) & \(5/7\) & \(16/21\) & 433\\
\(2++\) & \(20/21\) & \(3/7\) & 968\\
\(2+-\) & \(19/21\) & \(4/7\) & 790\\
\(2--\) & \(33/35\) & \(4/7\) & 1015\\
\bottomrule
\end{tabular}
\end{center}

Therefore:
\[
\boxed{q_*=1015.}
\]

The worst local pattern is \(2--\).

For the two-literal true branches, the corresponding exact local
characteristics are:
\[
\lambda_{2++}=\frac9{20},
\qquad
\lambda_{2+-}=\frac9{19},
\qquad
\lambda_{2--}=\frac5{11}.
\]

\section{The explicit high-incidence cutoff}

LINK-DICHOTOMY1 uses:
\[
B_4(q)=1,
\]
\[
B_s(q)
=
1+18(4-s)(q-1)B_{s+1}(q).
\]

Substitute \(q=1015\):
\[
B_4=1,
\]
\[
B_3=18253,
\]
\[
B_2=666307513,
\]
\[
B_1=36484334181829.
\]

Thus:
\[
\boxed{
D_*
=
18B_1
=
656718015272922.
}
\]

This number is intentionally interpreted structurally rather than
practically.

\begin{provedbox}
Within the declared resolved rank-\(\le4\) type-1-chain model:

\[
\Delta\le D_*
\Rightarrow
\text{fixed positive hypermatching coverage},
\]
while
\[
\Delta>D_*
\Rightarrow
\text{a favorable high-incidence local branch}.
\]

Therefore concentration is no longer an unstructured local escape in this
model.
\end{provedbox}

The bounded-incidence coverage guarantee is:
\[
\boxed{
f_{\rm covered,act}
\ge
\frac1{8D_*}
=
\frac1{5253744122183376}.
}
\]

The fraction is extremely small but independent of \(n\).

\section{What HIGHINC-CLOSE1 actually closes}

It closes:
\begin{itemize}[leftmargin=1.4em]
\item the finite one-/two-literal center-pattern audit;
\item the existence of one uniform finite petal threshold;
\item the local high-incidence DLS branch;
\item the bounded-incidence versus high-incidence structural dichotomy in the
declared resolved rank-\(\le4\) model.
\end{itemize}

It does not yet close:
\begin{itemize}[leftmargin=1.4em]
\item global integration into every stage of the deterministic solver;
\item every residual clause schema generated by broader preprocessing;
\item a new fully optimized global exact-3SAT exponent;
\item any lower bound for arbitrary SAT algorithms;
\item P versus NP.
\end{itemize}

This scope distinction is essential.

\section{Drive-recovered q-Horn/expander localization results}

The Drive corpus contains another genuine localization program based on proper
exact-2P2N formulas, q-Horn defect, and expansion.

The independently rechecked results include:
\begin{itemize}[leftmargin=1.4em]
\item additive q-Horn defect;
\item localization of defect to a majority residual incidence component;
\item a majority-component lineage theorem;
\item linear component-backdoor depth under the declared expansion hypotheses.
\end{itemize}

These theorems are useful because they show an explicit example of persistent
structural obstruction localizing inside a component over a linear
restriction prefix.

They do not solve CSE.

The same Drive reports retain:
\begin{itemize}[leftmargin=1.4em]
\item the affine/Tseitin bypass;
\item possible heterogeneous backdoor escapes;
\item representation sensitivity;
\item the missing arbitrary-decider acquisition/use theorem;
\item the absence of a nonrelativizing/nonalgebrizing transfer.
\end{itemize}

So this line remains a restricted structural control.

\section{Drive-recovered SFSL and orbit-spectral recurrence}

The live LSPV/RMV snapshot contains two independently verified exact lemmas.

\subsection{Seed-Fixed Surjectivity Lemma}

For distinct orbit points \(P_1,\ldots,P_m\) and degree budget at least
\(m-1\), low-degree interpolation can realize an arbitrary target vector on
those points.

This means generic low-degree source geometry is too expressive to force a
nonzero universal relation even on one fixed seed window.

That is an important negative result.

\subsection{Orbit-Spectral Recurrence Lemma}

For a fixed source polynomial \(p\), let:
\[
U_p
=
\operatorname{span}
\{p,T_Ap,T_A^2p,\ldots\},
\]
where:
\[
T_Af=f\circ A.
\]

If:
\[
s=\dim U_p,
\]
then the orbit evaluations satisfy a linear recurrence of order \(s\).

This creates a possible compact source-specific relation when \(s\) is small.

\subsection{The Booleanization gap}

The RMV distinguisher ultimately acts on Boolean output bits.

A field recurrence over an odd prime field does not automatically become a
small Boolean relation after projection through the actual binary code
coordinate.

So:
\[
\boxed{
\text{small field orbit dimension}
\not\Rightarrow
\text{small Boolean separator}
}
\]
without a new transfer theorem.

This remains open.

\section{How the exact-3SAT and SRAS tracks inform each other}

The two tracks operate at different logical levels.

The exact-3SAT track says:
\begin{quote}
Here are increasingly sophisticated transformations that really are cheap.
Any general hardness accounting must credit them.
\end{quote}

SRAS says:
\begin{quote}
Once all cheap transformations are credited, identify where the remaining
future-relevant distinction lives and charge only what ordinary computation
must actually pay.
\end{quote}

The relationship is therefore adversarial and useful.

Examples:

\begin{itemize}[leftmargin=1.4em]
\item Mediator elimination tells SRAS not to charge constant-width hidden
binary relations as hard.

\item Hypermatching tells SRAS that spread-out local correlations can be
aggregated efficiently.

\item High-incidence branching shows concentration need not imply hardness; it
may instead expose a useful localization/branching structure.

\item q-Horn and XOR controls show that large structural defect can coexist
with polynomial algorithms in another representation.

\item Replay results show that source history alone cannot be charged if the
input remains cheaply rereadable.
\end{itemize}

\section{The universal-evaluator branch}

Suppose CSE does not find a genuinely smaller source object.

Then the computation may instead maintain a compact global summary that
services all future restrictions.

For \(W_{\rm EXT}\), this is conceptually a reusable extension evaluator:
\[
E(F,\rho)
=
\operatorname{SAT}(F\upharpoonright\rho).
\]

Under the counterfactual \(P=NP\), the SAT algorithm itself is such an
evaluator.

Therefore a generic theorem saying:
\[
\text{universal evaluator is expensive}
\]
would already be extremely close to the original lower-bound problem.

This is why the project has moved away from trying to prove evaluator hardness
from:
\begin{itemize}[leftmargin=1.4em]
\item provenance;
\item explicit strategy-tree size;
\item fixed-level PH witness complexity;
\item random challenge freshness;
\item raw state multiplicity.
\end{itemize}

A successful theorem needs a more structural implementation-dependent reason.

\section{Wrapper covariance and equivalent implementations}

Any universal source-sensitive theorem must survive harmless changes to the
program.

The project repeatedly attacks candidate source properties with:
\[
\text{DEAD},\quad
\text{COPY},\quad
\text{DELAY},\quad
\text{INTERP},\quad
\text{RECOMP}.
\]

A raw property of syntax that changes under one of these wrappers cannot be a
universal necessary property of SAT algorithms.

But requiring complete implementation invariance is also too strong.  It can
collapse the property into a purely extensional function of the input/output
relation.

The preferred notion is covariance:
\[
E(w(A),F)
=
S_w(E(A,F)),
\]
where \(S_w\) is an explicitly controlled transport map whose cost is charged.

This allows implementation-sensitive information without pretending different
implementations are literally identical.

\section{Why easy-family controls are mandatory}

Any proposed hardness resource is tested against known polynomial-time
families.

\subsection{2-SAT}

Strong graph connectivity can summarize global implication structure.

\subsection{XOR-SAT / Tseitin}

Gaussian elimination can compress apparently global combinatorial dependence
into a linear basis.

\subsection{Horn and dual-Horn}

Closure procedures can propagate information globally in polynomial time.

\subsection{Bounded treewidth}

Dynamic programming across separators can solve formulas with huge assignment
spaces but small structural width.

\subsection{Independent products}

Semantic multiplicity can grow exponentially while the decision problem
factorizes trivially.

If a candidate resource becomes large on one of these families, the resource
may still be interesting, but it is not by itself an unrestricted time measure.

\section{Retired proof mechanisms}

The following generic implications have been retired.

\begin{longtable}{p{0.40\textwidth} p{0.52\textwidth}}
\toprule
\textbf{Claim} & \textbf{Reason it fails as a generic unrestricted step}\\
\midrule
\endhead
Many continuations imply long runtime &
They imply distinguishability, not acquisition time.\\
Large explicit representation implies long runtime &
Another representation may be succinct.\\
Unique provenance implies expensive recovery &
Static reread or deterministic rerun may be cheap.\\
Fresh random challenge implies expensive answer &
A universal evaluator can answer after the challenge appears.\\
Long constant-width inference implies hard discovery &
Constant-width closure is polynomially enumerable.\\
Large OBDD/MDD/tree representation implies runtime &
Representation-specific lower bounds need an extraction theorem.\\
Large rank implies runtime &
Algebraic bases can be compact and efficiently computable.\\
Low-degree mediator structure implies easy SAT &
The residual clean-core schema remains NP-complete.\\
Low outside degree implies large matching &
Relations can concentrate on a common chain.\\
Concentration implies hardness &
HIGHINC-CLOSE1 converts sufficiently high concentration into a favorable
local branch in the declared model.\\
\bottomrule
\end{longtable}

\section{What would constitute an actual P-vs-NP proof in this framework}

A completed proof would need all of the following with explicit quantifiers.

\subsection{Step A: arbitrary algorithm}

Fix an arbitrary deterministic polynomial-time exact SAT decider \(A\).

The proof cannot assume \(A\) resembles DPLL, resolution, an OBDD, or any
known solver architecture.

\subsection{Step B: forced SAT-specific obligation}

Show that correctness on a hard family forces \(A\) to support enough
future-relevant behavior to amount to:
\[
n-O(\log n)
\]
valid reduction in the independently defined consumer model.

This step is currently open.

\subsection{Step C: source/accounting classification}

Apply SRAS-ROUTECLOSE1.

The required distinction must be:
\[
\text{retained/localized}
\vee
\text{reacquirable}
\vee
\text{external}.
\]

For ordinary computation, EXTERNAL cannot simply be assumed available.

\subsection{Step D: quantitative ordinary-source theorem}

Prove:
\[
\mathcal C_{F_n,U_n,M}
\bigl(n-O(\log n)\bigr)
=
n^{\omega(1)}
\]
for every legal retained/reacquired ordinary realization.

This is the central missing lower bound.

\subsection{Step E: barrier discharge}

Identify the load-bearing line that fails under arbitrary oracle replacement.

ASCS supplies a plausible source-semantic candidate.

Then separately show the proof does not merely algebrize.

\subsection{Step F: exact model-matched contradiction}

Only after all prior arrows hold may the recursive closure/time-space endpoint
be invoked.

\section{Quantifier discipline}

The proof goal is not:
\[
\exists F_n\ \forall A
\]
unless a universal hard family is actually proved.

A more natural lower-bound quantifier shape is:
\[
\forall A\in\mathrm{DTIME}(n^{O(1)})
\ \exists^\infty n\
\exists F\in 3\mathrm{SAT}_n
\]
such that \(A\) exceeds the alleged polynomial bound or fails.

Every extraction theorem must state which quantities depend on \(A\), which
depend on \(F\), and which are fixed in advance.

This discipline prevents a common hidden inversion of quantifiers.

\section{Current dependency graph}

The current state can be summarized as:

\[
\boxed{
\begin{array}{rcl}
\text{BRTHRESH1} &:& \text{proved accounting}\\
\text{CONTSEP1/NO2K1} &:& \text{proved representation constraints}\\
\text{SRAS-ROUTECLOSE1} &:& \text{proved route completeness}\\
\text{EXTCONSUMER1} &:& \text{proved fixed consumer}\\
\text{TRACECLOSE1/FRESH-EVAL1/WIDTHCLOSE1} &:& \text{proved no-go controls}\\
\text{RECURCLOSE1/MULTICHANNEL1} &:& \text{proved conditional closure}\\
\text{CSE1} &:& \text{open}\\
\text{near-total SAT reduction forcing} &:& \text{open}\\
\text{ordinary-source acquisition LB} &:& \text{open}\\
\text{ASCS oracle break} &:& \text{proved diagnostic}\\
\text{FSC} &:& \text{open}\\
\text{nonalgebrizing hinge} &:& \text{open}\\
\text{P}\ne\text{NP} &:& \text{open}.
\end{array}}
\]

On the exact-3SAT side:
\[
\boxed{
\begin{array}{rcl}
\text{mediator elimination} &:& \text{proved}\\
\text{resolved characteristic formula} &:& \text{proved}\\
\text{rank-2/3/4 finite classification} &:& \text{proved}\\
\text{active density} &:& \text{proved}\\
\text{hypermatching density} &:& \text{proved}\\
\text{link dichotomy} &:& \text{proved}\\
\text{five center-pattern branch audit} &:& \text{proved}\\
\text{resolved-incidence dichotomy} &:& \text{proved in stated model}\\
\text{full pipeline integration} &:& \text{open}.
\end{array}}
\]

\section{Research priorities after the current synthesis}

\subsection{Priority 1: integrate HIGHINC-CLOSE1 globally}

The local concentration escape is closed in the declared resolved model.

The next exact-algorithm task is to verify that every residual clause form
created by the full preprocessing pipeline falls under the proved accounting,
or else to add the missing case.

Then compute one global deterministic balance.

\subsection{Priority 2: sharpen the SAT-specific extraction statement}

Use \(W_{\rm EXT}\) as the fixed consumer.

Try to prove:
\[
\boxed{
\text{fractional source localization}
\vee
\text{globally supported reusable evaluator}.
}
\]

This would be a more precise CSE dichotomy than asking an arbitrary trace to
look locally decomposable.

\subsection{Priority 3: attack FSC inside the exact source map}

Do not ask for a generic certificate for arbitrary long computation.

Work inside the actual canonical proof
\[
\to
\text{low-degree extension}
\to
\text{RMV output}
\]
map.

Determine the cheapest exact ordinary certificate for:
\[
\text{selector }\sigma\text{ produces output }y.
\]

Measure both witness length and deterministic verification circuit size.

\subsection{Priority 4: Booleanization of source-specific recurrence}

If an accepted proof polynomial has small orbit-spectral dimension, determine
whether the exact binary conversion preserves any compact Boolean relation.

A negative theorem would retire the spectral separator route cleanly.

A positive theorem could feed the compact-separator/reconstruction branch.

\subsection{Priority 5: nonalgebrization audit}

Any successful source theorem must be tested under algebraic-oracle extension.

Do not infer nonalgebrization merely from nonrelativization.

\section{Source basis for this exposition}

This document is based on the current project records, especially:

\begin{itemize}[leftmargin=1.4em]
\item \texttt{main2.tex};
\item \texttt{research/explorations/source-reduction/2026-10-05\_SRAS\_CSE1.md};
\item \texttt{research/explorations/source-reduction/2026-10-05\_HIGHINC\_CLOSE1.md};
\item the high-incidence experiment and result files;
\item the random/replay and recursive-closure GitHub records listed in
\texttt{main2.tex};
\item the Drive LSPV source manual containing ASCS/FSC;
\item the Drive MAIN LSPV/RMV snapshot containing SFSL and the orbit-spectral
recurrence;
\item the independently verified Drive q-Horn/expander manuscript chain.
\end{itemize}

The Drive-derived restricted results are included here only at the scope
supported by those records.  In particular, component-backdoor depth is not
promoted to unrestricted SAT runtime, field recurrences are not promoted to
Boolean separators, and the ASCS oracle-break diagnostic is not promoted to a
P-vs-NP lower bound.

\section{Final perspective}

The program has made progress mostly by removing false shortcuts.

It now has:
\begin{itemize}[leftmargin=1.4em]
\item a disciplined semantic reduction measure;
\item a route-complete accounting framework;
\item a fixed solver-independent continuation consumer;
\item exact recursive-closure formulas;
\item strong replay, provenance, width, and representation no-go controls;
\item a concrete ordinary-source predicate that genuinely can fail
relativization;
\item an increasingly complete exact-3SAT structural laboratory;
\item a closed local high-incidence case in the resolved rank-\(\le4\)
type-1-chain model.
\end{itemize}

What it still does not have is the theorem that converts arbitrary polynomial
SAT computation into an ordinary-source object whose required acquisition or
maintenance cost is superpolynomial.

That is the remaining central arrow:
\[
\boxed{
\text{arbitrary polynomial SAT computation}
\Longrightarrow
\text{forced source-matched contraction or superpolynomial acquisition}.
}
\]

Until that arrow is proved, together with the necessary barrier discharge:
\[
\boxed{\text{$P$ versus $NP$ remains OPEN.}}
\]

\appendix

\section{Appendix A: concise glossary}

\begin{longtable}{p{0.22\textwidth} p{0.69\textwidth}}
\toprule
\textbf{Term} & \textbf{Meaning}\\
\midrule
\endhead
SRAS & Source-Reduction Accounting System.\\
CSE & Contractive Source Extraction.\\
\same & Require the object to belong to the same committed source realization.\\
Continuation equivalence & Indistinguishability under every declared legal future consumer.\\
Acquisition curve & Minimum charged cost needed to obtain a declared amount of valid reduction.\\
Universal evaluator & Reusable representation/algorithm answering every legal future challenge.\\
Reacquisition & Reread, rerun, replay, or recomputation of discarded information.\\
ASCS & Absolute Source-Complexity Separator.\\
FSC & Fast Source Certificate needed to make ASCS compatible with the small reconstruction interface.\\
Type-1 chain & Positive 3-clause with seven satisfying local states.\\
Resolved relation & Finite-domain relation exposed after mediator elimination.\\
Hypermatching & Pairwise chain-disjoint set of resolved relations.\\
Link hypergraph & Outside supports of relations incident to a common high-incidence chain.\\
Petal & Portion of a resolved support outside the common center.\\
HIGHINC-CLOSE1 & Finite five-pattern high-incidence audit and resulting structural cutoff.\\
\bottomrule
\end{longtable}

\section{Appendix B: status legend}

\begin{description}[leftmargin=1.8em,style=nextline]
\item[PROVED/EXACT]
Proved for the explicitly declared model or exhaustively verified finite
statement with an exact mathematical derivation.

\item[CONDITIONAL]
The implication is proved assuming an explicitly named open hypothesis.

\item[CANDIDATE]
A proposed theorem, object, or proof architecture whose load-bearing content is
not yet proved.

\item[COMPUTATIONAL]
Finite experiment or numerical check only.

\item[FALSIFIED/RETIRED]
A generic implication has a counterexample, closure theorem, wrapper escape, or
other reason it cannot carry the intended unrestricted conclusion.
\end{description}

\section{Appendix C: theorem-to-gap map}

\begin{longtable}{p{0.26\textwidth} p{0.30\textwidth} p{0.34\textwidth}}
\toprule
\textbf{Result} & \textbf{What it gives} & \textbf{What it does not give}\\
\midrule
\endhead
BRTHRESH1 & Required reduction scale & Arbitrary-algorithm extraction\\
CONTSEP1 & Safe-summary range lower bound & Runtime lower bound\\
SRAS-ROUTECLOSE1 & Exhaustive location of future distinctions & Cost of ordinary routes\\
EXTCONSUMER1 & Fixed SAT continuation semantics & Fractional source shrinkage\\
TRACECLOSE1 & Trace closure no-go & Structural CSE\\
FRESH-EVAL1 & Freshness no-go & Evaluator hardness\\
WIDTHCLOSE1 & Fixed-width closure is cheap & Higher-level acquisition LB\\
RECURCLOSE1 & Exponent after valid contraction & Existence of contraction\\
ASCS & Concrete oracle-break source predicate & Small exact verification/FSC\\
CLAUSECHAR1 & Exact resolved relation characteristic & Global solver integration\\
LINK-DICHOTOMY1 & High incidence gives disjoint petals or deeper center & Branch gain by itself\\
LOCALPATTERN-BRANCH1 & Branch gain for all five local sign types & Full preprocessing coverage\\
RESOLVED-INCIDENCE-DICHOTOMY1 & Bounded/high incidence both handled locally & Unrestricted SAT lower bound\\
\bottomrule
\end{longtable}

\begin{thebibliography}{9}

\bibitem{Liu2018}
S. Cliff Liu,
\emph{Chain, Generalization of Covering Code, and Deterministic Algorithm for $k$-SAT},
ICALP 2018 / arXiv:1804.07901.

\bibitem{BGS1975}
T. Baker, J. Gill, R. Solovay,
\emph{Relativizations of the $P=?NP$ Question},
SIAM Journal on Computing 4(4), 1975.

\bibitem{AW2009}
S. Aaronson, A. Wigderson,
\emph{Algebrization: A New Barrier in Complexity Theory},
ACM Transactions on Computation Theory, 2009.

\bibitem{RR1997}
A. Razborov, S. Rudich,
\emph{Natural Proofs},
Journal of Computer and System Sciences 55(1), 1997.

\bibitem{BW2001}
E. Ben-Sasson, A. Wigderson,
\emph{Short Proofs are Narrow---Resolution Made Simple},
Journal of the ACM 48(2), 2001.

\bibitem{IPZ2001}
R. Impagliazzo, R. Paturi, F. Zane,
\emph{Which Problems Have Strongly Exponential Complexity?},
Journal of Computer and System Sciences 63(4), 2001.

\end{thebibliography}

\end{document}
